\documentclass[12pt, letterpaper]{article}

% =========================================================
% PAQUETES DE CONFIGURACIÓN Y IDIOMA
% =========================================================
\usepackage[spanish, es-tabla]{babel}
\usepackage[utf8]{utf8}
\usepackage[margin=2.5cm]{geometry}
\usepackage{amsmath, amssymb}
\usepackage{graphicx}
\usepackage{float}
\usepackage{array}
\usepackage{booktabs}
\usepackage{xcolor}
\usepackage{hyperref}

\hypersetup{
    colorlinks=true,
    linkcolor=blue,
    filecolor=magenta,      
    urlcolor=cyan,
    citecolor=blue,
    pdftitle={Reporte de Métricas y Calidad - Proyecto Outline},
    pdfpagemode=FullScreen,
}

% =========================================================
% ENCABEZADO Y DATOS DEL PROYECTO
% =========================================================
\title{\textbf{Proyecto Integrador: Outline}\\ \large Subsección 4.4: Investigación de Métricas Reportables e Interpretación}
\author{\textbf{Noé Ezequiel Estrada Ramirez} \\ Líder / Analista --- Equipo: \texttt{TIID05B\_Equipo4} \\ Docente: Dr. Raul Velasquez}
\date{\today}

\begin{document}

\maketitle

\hrule
\vspace{0.5cm}

% =========================================================
% SECCIÓN 4.4: INVESTIGACIÓN DE MÉTRICAS REPORTABLES E INTERPRETACIONAL
% =========================================================

\section{Estrategia y Métricas de Calidad de Software}

\subsection{Investigación: métricas reportables e interpretación}

La evaluación cuantitativa del software requiere la definición formal de métricas tanto de producto (código fuente y estructura) como de proceso de pruebas \cite{iso25010, galin2018}. A continuación, se clasifican las métricas clave, sus umbrales de referencia y su correspondencia con el modelo de calidad ISO/IEC 25010:2023 e ISO/IEC 5055:2021 \cite{iso25023, iso5055}.

\begin{table}[H]
\centering
\small
\caption{Métricas reportables de código y pruebas ligadas a ISO/IEC 25010:2023.}
\label{tab:metricas_calidad}
\begin{tabular}{|p{2.5cm}|p{3.2cm}|p{3.5cm}|p{2.8cm}|p{3.0cm}|}
\hline
\textbf{Categoría} & \textbf{Métrica} & \textbf{Fórmula / Indicador} & \textbf{Umbral de Referencia} & \textbf{Característica ISO 25010} \\ \hline
Código & Complejidad Ciclomática $V(G)$ & $V(G) = P + 1$ \cite{mccabe1976} & $V(G) \le 10$ por función & Mantenibilidad (Analizabilidad) \\ \hline
Código & Halstead (Volumen $V$) & $V = (N_1 + N_2) \log_2 (n_1 + n_2)$ & $V \le 1000$ bits/módulo & Mantenibilidad (Comprensibilidad) \\ \hline
Código & Puntos de Función (IFPUG) & Conteo funcional de ILF, EIF, EI, EO, EQ & Según la historia ($\le 5$ FP/HU) & Adecuación Funcional \\ \hline
Estructural & Reglas ISO/IEC 5055 & Violaciones estáticas en 4 pilares & 0 fallas críticas & Seguridad y Mantenibilidad \\ \hline
Pruebas & Cobertura de Código & $\frac{\text{Líneas Ejecutadas}}{\text{Líneas Totales}} \times 100$ & $\ge 80\%$ (Outline: $62\%$) & Testabilidad (Mantenibilidad) \\ \hline
Pruebas & Densidad de Defectos & $\frac{\text{Defectos Conf.}}{\text{KLOC}}$ & $\le 1.0$ defectos/KLOC & Confiabilidad (Madurez) \\ \hline
Pruebas & Tasa de Aprobación & $\frac{\text{Casos Aprobados}}{\text{Casos Ejecutados}} \times 100$ & $\ge 95\%$ (Outline: $100\%$) & Correctitud Funcional \\ \hline
\end{tabular}
\end{table}

\subsubsection{Análisis e Interpretación de Métricas de Código y Pruebas}

\begin{itemize}
    \item \textbf{Complejidad Ciclomática de McCabe ($V(G)$):} Determina el número de caminos independientes en el flujo de ejecución \cite{mccabe1976}. Valores superiores a 10 reflejan baja mantenibilidad y exigen división modular.
    \item \textbf{Métricas de Halstead:} Evalúan la complejidad cuantitativa midiendo operadores ($n_1, N_1$) y operandos ($n_2, N_2$) únicos y totales \cite{galin2018}. Un volumen o dificultad elevado incrementa la probabilidad de inconsistencias lógicas.
    \item \textbf{Métricas Estructurales ISO/IEC 5055:2021:} Analizan cuantitativamente las vulnerabilidades en código fuente que impactan en la confiabilidad y mantenibilidad, detectando patrones indeseables antes del despliegue \cite{iso5055}.
    \item \textbf{Métricas de Pruebas (Cobertura y Densidad):} La cobertura de código permite verificar la efectividad de los casos de prueba ejecutados (e.g. Selenium/PyTest). La densidad de defectos valida la estabilidad de las funcionalidades del sistema \cite{istqb2024}.
\end{itemize}

\subsubsection{Cuadro de KPIs Definidos para el Proyecto Outline}

A partir de la especificación de calidad de la wiki interna \textit{Outline}, se consolidan los 7 KPIs clave de desempeño y calidad de software \cite{pressman2020}:

\begin{enumerate}
    \item \textbf{KPI-01 (Complejidad Ciclomática):} $V(G) \le 8$ promedio por módulo de edición.
    \item \textbf{KPI-02 (Duplicación de Código):} Duplicación $< 3\%$ según reporte estático SonarQube / ISO 5055.
    \item \textbf{KPI-03 (Cobertura de Pruebas):} Cobertura de código automatizado $\ge 80\%$ (Meta de refactorización respecto al $62\%$ de P2).
    \item \textbf{KPI-04 (Densidad de Defectos):} $< 0.8$ defectos por KLOC en etapa de integración.
    \item \textbf{KPI-05 (Tasa de Aprobación de Pruebas):} $100\%$ de aprobación en la suite de regresión funcional (42/42 pruebas pasadas).
    \item \textbf{KPI-06 (Seguridad en Código):} 0 vulnerabilidades de severidad Alta/Bloqueante.
    \item \textbf{KPI-07 (Comportamiento Temporal):} Tiempo de respuesta $< 500$\,ms al renderizar notas Markdown.
\end{enumerate}

\vspace{0.8cm}
\hrule
\vspace{0.5cm}

% =========================================================
% SECCIÓN 6: APORTE INDIVIDUAL
% =========================================================

\section{Aportes Individuales}

\subsection*{Aporte Individual: Noé Ezequiel Estrada Ramirez (Líder / Analista)}

Durante el desarrollo de esta investigación, coordiné la estructuración analítica de las métricas de calidad de software aplicables al proyecto Outline, integrando las especificaciones de ISO/IEC 25010:2023, ISO/IEC 25023 e ISO/IEC 5055:2021. Fundamenté la relación entre las métricas de código (McCabe, Halstead e IFPUG) y las métricas de pruebas automáticas (cobertura y densidad de defectos), definiendo los umbrales de tolerancia técnica y consolidando el conjunto de KPIs del proyecto. Este análisis permitió establecer criterios claros para interpretar los resultados del análisis estático y dinámico en el reporte integrador.

\vspace{0.8cm}

% =========================================================
% BIBLIOGRAFÍA DE REFERENCIA
% =========================================================

\begin{thebibliography}{9}

\bibitem{iso25010}
ISO/IEC 25010:2023, \textit{Systems and software engineering --- Systems and software Quality Requirements and Evaluation (SQuaRE) --- Product quality model}, International Organization for Standardization, 2023.

\bibitem{iso25023}
ISO/IEC 25023:2016, \textit{Systems and software engineering --- SQuaRE --- Measurement of system and software product quality}, International Organization for Standardization, 2016.

\bibitem{iso5055}
ISO/IEC 5055:2021, \textit{Software engineering --- Software product quality measurement --- Automated source code quality metrics}, International Organization for Standardization, 2021.

\bibitem{mccabe1976}
McCabe, T. J. (1976). A Complexity Measure. \textit{IEEE Transactions on Software Engineering}, SE-2(4), 308-320.

\bibitem{galin2018}
Galin, D. (2018). \textit{Software Quality Assurance: From theory to implementation}. John Wiley \& Sons.

\bibitem{istqb2024}
ISTQB. (2024). \textit{Certified Tester Foundation Level (CTFL) Syllabus v4.0}. International Software Testing Qualifications Board.

\bibitem{pressman2020}
Pressman, R. S., \& Maxim, B. R. (2020). \textit{Software Engineering: A Practitioner's Approach} (9th ed.). McGraw-Hill Education.

\end{thebibliography}

\end{document}